\documentclass[../../../include/open-logic-section]{subfiles}

\begin{document}

\olfileid{sth}{choice}{woproblem}
\olsection{Prakara Pengurutan Becik}

Cetha yèn akèh kang gumantung marang apa kita nampa Pengurutan Becik.
Nanging rembugan bab prinsip iki lumrahé tumuju marang prinsip kang
ekuivalen, yaiku Aksioma Pilihan. Mula saiki kita nggatekaké aksioma
kasebut lan mbuktekaké ekuivalensiné.

Ing \citeyear{Cantor1883}, Cantor nyengkuyung Aksioma Pengurutan Becik,
lan nyebut iku ``hukum pikiran kang katon dhasar, sugih akibaté,
lan mligi banget amarga lumakune sacara umum'' (dipetik ing
\citeauthor{Potter2004} \citeyear[kaca~243]{Potter2004}). Nanging
pungkasané Cantor yakin yèn ``Aksioma'' iku perlu bukti. Semono uga
komunitas matématika.

Prakara iki ``dipecahaké'' déning Zermelo ing
\citeyear{Zermelo1904}. Kanggo nerangaké cara pamungkasé, kita
mbutuhaké sawatara dhéfinisi.

\begin{defn}
Fungsi $f$ diarani \emph{fungsi pilihan} yèn lan mung yèn
$f(x) \in x$ kanggo saben $x \in \dom{f}$. Kita nyebut $f$
\emph{fungsi pilihan kanggo $A$} yèn $f$ fungsi pilihan kanthi
$\dom{f} = A \setminus \{\emptyset\}$.
\end{defn}

Miturut pangertèn intuitif, kanggo saben himpunan kang ora kosong
$x \in A$, fungsi pilihan kanggo $A$ \emph{milih} siji unsur tartamtu,
$f(x)$, saka $x$. Aksioma Pilihan banjur mratelakaké:

\begin{axiom}[Pilihan]
	Saben himpunan nduwèni fungsi pilihan.
\end{axiom}

Zermelo nuduhaké yèn Pilihan nyebabaké Pengurutan Becik, lan kosok baliné:

\begin{thm}[ing $\ZF$]\ollabel{thmwochoice}
Pengurutan Becik lan Pilihan iku ekuivalen.
\end{thm}

\begin{proof}
\emph{Saka kiwa menyang tengen.} Jupuken $A$ minangka himpunan saka
himpunan. Mula $\bigcup A$ ana miturut Aksioma Gabungan, saéngga
miturut Pengurutan Becik ana $<$ kang ngurutaké $\bigcup A$ kanthi becik.
Saiki tetepna $f(x) = \text{unsur paling cilik miturut $<$ saka }x$.
Iki fungsi pilihan kanggo $A$.

\emph{Saka tengen menyang kiwa.} Tetepna $A$. Yèn himpunan iki kosong,
dhèwèké wis terurut becik; mula anggepen ora kosong. Miturut Pilihan,
ana fungsi pilihan $f$ kanggo $\Pow{A} \setminus \{\emptyset\}$.
Kanthi Rekursi Transfinit, dhéfinisikaké fungsi:
\begin{align*}
	g(0) &= f(A)\\
	g(\alpha) &=
		\begin{cases}
			\text{mandheg!{}} &\text{yèn }A \subseteq \funimage{g}{\alpha}\\
			f(A \setminus \funimage{g}{\alpha}) & \text{yèn ora}\\	
		\end{cases}
\end{align*}
Pratandha ``mandheg!'' iku cekakan saka dhéfinisi kang luwih dawa:
nalika $A = \funimage{g}{\alpha}$ pisanan kelakon, tetepna
$g(\delta) = A$ kanggo saben $\delta \geq \alpha$ minangka nilai
panandha, banjur gunakna mung watesan sadurungé tahap iku. Ing wiwitan
kita mung ngerti yèn rekursi iki ndhéfinisikaké sawijining \emph{suku}
(delengen cathetan sawisé
\olref[sth][spine][recursion]{transrecursionschema}); ing ngisor iki
kita mbuktekaké yèn watesan kang dibutuhaké iku sawijining fungsi.

Amarga $f$ fungsi pilihan, kanggo saben $\alpha$ sadurungé mandheg
kita nduwèni $g(\alpha) =
f(A \setminus \funimage{g}{\alpha}) \in A \setminus
\funimage{g}{\alpha}$; yaiku $g(\alpha) \notin \funimage{g}{\alpha}$.
Mula yèn $g(\alpha) = g(\beta)$, banjur $g(\beta) \notin
\funimage{g}{\alpha}$, yaiku $\beta \notin \alpha$; semono uga
$\alpha \notin \beta$. Miturut Trikotomi, $\alpha = \beta$.
Dadi watesan $g$ sadurungé mandheg iku !!{injective}.

Sabanjuré, proses iki mesthi mandheg, yaiku ana ordinal $\alpha$ paling
cilik kang nyukupi $A = g[\alpha]$. Yèn ora, watesan $g$ bakal
!!{injective} lan kita bakal nduwèni $\cardle{\alpha}{A}$ kanggo
saben ordinal $\alpha$, nalisir \olref[hartogs]{HartogsLemma}.
Sawisé fungsi iki diwatesi ing segmen sadurungé mandheg, kita uga nduwèni
$\ran{g} = A$.

Saka kabèh iki, watesan $g$ iku !!a{bijection} saka sawijining ordinal
menyang $A$. Saiki $g$ bisa digunakaké kanggo ngurutaké $A$ kanthi becik.
\end{proof}

Dadi Pengurutan Becik lan Pilihan padha-padha bener utawa padha-padha
gagal. Nanging pitakonané isih ana: apa kaloroné bener?

\end{document}
